% from main.tex
%%%%%%%%%%  scan idx 160  |  folio 147  %%%%%%%%%%
\emph{Preuve.} Il suffit de
montrer que l'homomorphisme induit par le rang sur le faisceau associé à \(k(\Dpx^-(S)_{\mathrm{parf}})\) est un homomorphisme de faisceaux d'anneaux compatible avec les images inverses. D'après (6.9), on peut remplacer \(k(\Dpx^-(S)_{\mathrm{parf}})\) par \(k(\mathcal C_1)\). On est donc ramené à vérifier a) et b) pour \(L\) et \(M\) libres de type fini, ce qui est trivial.

\textbf{6.13.} \((S,A)\) désignant toujours un topos localement annelé, soient \(X\in\Ob S\), et \(\underline n\in H^0(X,\mathbb Z_S)\). On définit un objet de \(\Parf(X)\), noté \(A_X^{\underline n}\), de la manière suivante. Si \(\underline n\) est constante de valeur un entier \(n\geq0\) (resp. \(n<0\)), on pose \(A_X^{\underline n}=A_X^n\) (resp. \(A_X^{\underline n}=A_X^{-n}[1]\)); dans le cas général, on décompose \(X\) en morceaux disjoints \(X_i\) tels que \(\underline n|X_i\) soit constante, et \(A_X^{\underline n}\) est défini comme l'unique objet de \(\Parf(X)\) induisant \(A_{X_i}^{\underline n|X_i}\) sur chaque \(X_i\). (On peut, plus généralement, définir \(L^{\underline n}\) pour \(L\in\Ob\Parf(X)\), en posant \(L^{\underline n}=A_X^{\underline n}\otimes L\).).

L'application \(i_X:H^0(X,\mathbb Z_S)\longrightarrow k(\Parf(X))\), définie par \(i_X(\underline n)=\cl(A_X^{\underline n})\), est un homomorphisme d'anneaux (tel que \(i_X(1)=1\)), et définit, pour \(X\) variable, un homomorphisme de préfaisceaux d'anneaux
\[
  i:\mathbb Z_S\longrightarrow k(\Parf(S)).
\]

\textbf{Proposition 6.14.} Sous les hypothèses de (6.13), on a
\[
  \rg\, i=\Id(\mathbb Z_S)
\]
(autrement dit, \(k(\Parf(S))\), considéré comme préfaisceau de \(\mathbb Z_S\)-algèbres au moyen de \(i\), est augmenté vers \(\mathbb Z_S\) par le rang). En outre le rang
\[
  \rg:k(\Parf(S))\longrightarrow\mathbb Z_S
\]
définit un isomorphisme sur les faisceaux associés.

